%! Symmetric Positive Definite Matrices — Unit 6, Lesson 6.3 — Exit Ticket (Answer Key)
\documentclass[10pt]{article}
\usepackage{linalg-article}
\usepackage{linalg-key}   % pulls in -boxes; do NOT also load -boxes
\newcommand{\TallMath}[1]{$\displaystyle #1\rule[-1.4em]{0pt}{3.2em}$}
\newcommand{\T}{\mathsf{T}}

\begin{document}

\pageheader{Unit 6, Lesson 6.3}{Exit Ticket --- Answer Key}
\namedateperiod

\vspace{0.15in}

No notes. Show your reasoning. Use $S=\begin{bmatrix} 4 & 1 \\ 1 & 4 \end{bmatrix}$, whose eigenvalues are
$5$ and $3$ with eigenvectors $\begin{bsmallmatrix}1\\1\end{bsmallmatrix}$ and $\begin{bsmallmatrix}1\\-1\end{bsmallmatrix}$.

\begin{enumerate}[leftmargin=1.8em, itemsep=14pt]
  \item \textbf{Symmetric \& perpendicular.} Explain why $S$ is symmetric, then use a dot product to check
        that its two eigenvectors are perpendicular.
        \ansline{$S=S^{\T}$: flipping across the diagonal keeps the off-diagonal $1$'s in place, so $S$ is symmetric. Dot product: $\begin{bsmallmatrix}1\\1\end{bsmallmatrix}\cdot\begin{bsmallmatrix}1\\-1\end{bsmallmatrix}=1-1=0$, so the eigenvectors are perpendicular.}
  \item \textbf{Positive definite?} Decide whether $S$ is positive definite \emph{two} ways: from the
        eigenvalues, and from the quick test ($a>0$ and $ac-b^2>0$).
        \ansline{Eigenvalues $5$ and $3$ are both positive $\Rightarrow$ positive definite. Quick test: $a=4>0$ and $ac-b^2=(4)(4)-1^2=15>0$ $\Rightarrow$ positive definite. Both agree.}
  \item \textbf{Explain.} How could you have known --- \emph{without} computing them --- that this matrix's
        eigenvectors would be perpendicular? (One or two sentences.)
        \ansline{Because $S$ is symmetric. The spectral theorem guarantees that every symmetric matrix has real eigenvalues and perpendicular eigenvectors --- no computation needed.}
\end{enumerate}

\begin{teachernote}
Sort into ``got it / arithmetic slip / concept slip.'' Item~1 checks recognizing symmetry and the
dot-product perpendicularity test; item~2 checks \emph{both} positive-definite tests (all $\lambda>0$; and
$a>0$ \& $ac-b^2=15>0$); item~3 is the target idea --- symmetry \emph{alone} forces perpendicular
eigenvectors (the spectral theorem). Watch for students who test positive definiteness from the diagonal
only, or who think perpendicular eigenvectors need to be computed to be known. If many miss item~3,
re-anchor on the §1 spectral-theorem statement before Lesson~6.4.
\end{teachernote}

\end{document}
